\chapter{Mode of Production and Challenges of Social Transformation}

\section{Dialectics and the Mode of Production Debate}

\begin{KeyConcept}
\begin{itemize}
\item \textbf{Dialectics} (Hegel and Marx) is the process in which \textbf{two contradictory ideas come into conflict and synthesise into a new idea}. Hegel's was \textbf{dialectical idealism}; Marx's was \textbf{dialectical materialism} (it is \textbf{being that determines consciousness}, not the reverse --- his real assumption).
\item The contradiction comes from the \textbf{labour process} (labour is the expression of human potential and also transforms humans); capitalism benefits one class and exploits another at the same time --- so \textbf{two antagonistic ideas coexist}, and their churning produces a new idea (communism is the culmination where the dialectic ends). Alienated workers rebuild community at work through \textbf{Elton Mayo's 'informal organisations'}.
\item \textbf{Paul Brass's 'dual modernity'}: whenever a new idea emerges, the indigenous/old idea \textbf{revives at the same time} (English education $\rightarrow$ vernacular revival; allopathic medicine $\rightarrow$ Ayurveda).
\item \textbf{Mode of production} = \textbf{means of production + relations of production} (and their inherent contradictions); Marx held there is one dominant mode at a time. The Indian \textbf{mode-of-production debate} asks whether, after Green Revolution, agriculture is capitalist, feudal, or semi-feudal/pre-capitalist.
\item \textbf{Weber vs Marx}: Weber is ideational (the \emph{spirit} of capitalism develops), Marx material; \textbf{Kautsky} (and Tawney) argued it was \textbf{capitalism that produced Protestantism}, not the reverse (capitalism already existed in Holland/Scandinavia before Protestantism).
\item \textbf{E. P. Thompson}: the agrarian--industrial shift made \textbf{time repressive} (in agriculture one worked at one's comfort; in industry the shift-hooter governs time).
\end{itemize}
\end{KeyConcept}

\begin{QuickRevision}
\begin{itemize}
\item Dialectics
\item two contradictory ideas come into conflict and synthesise into a new idea
\item dialectical idealism
\item dialectical materialism
\item being that determines consciousness
\item labour process
\item two antagonistic ideas coexist
\item Elton Mayo's 'informal organisations'
\item Paul Brass's 'dual modernity'
\item revives at the same time
\item Mode of production
\item means of production + relations of production
\item mode-of-production debate
\item Weber vs Marx
\end{itemize}
\end{QuickRevision}

\section{The Subaltern School}

\begin{KeyConcept}
\begin{itemize}
\item One of the \textbf{two great contributions of Indian social science} (the other being the mode-of-production debate) is the \textbf{subaltern school} --- history/sociology \textbf{from the perspective of the marginalized}.
\item It dismantles the mainstream idea that only the 'enlightened' have true consciousness. Example: tribals fighting the British with primitive weapons were called 'irrational' --- but for them, whose subsistence was destroyed by forest and agricultural policies, fighting was the \textbf{most rational choice} (death was imminent anyway).
\item \textbf{Dalit sociology vs subaltern sociology}: Dalit sociology can be written from an elite perspective (Mulk Raj Anand's \emph{Untouchable}, Premchand --- where a Dalit is emancipated only with upper-caste help); subaltern sociology writes from the \textbf{consciousness of the Dalits themselves}. Example: \textbf{Kancha Ilaiah} (English-medium education as an instrument of Dalit emancipation) vs the Dalit activist who argued that vernacular education keeps Dalits rooted in their community (English-educated Dalits abandon it) --- two different perspectives on the same empirical reality.
\end{itemize}
\end{KeyConcept}

\begin{QuickRevision}
\begin{itemize}
\item two great contributions of Indian social science
\item subaltern school
\item from the perspective of the marginalized
\item most rational choice
\item Dalit sociology vs subaltern sociology
\item consciousness of the Dalits themselves
\item Kancha Ilaiah
\end{itemize}
\end{QuickRevision}

\section{Amit Bhaduri and Utsa Patnaik: Semi-Feudal vs Pre-Capitalist}

\begin{KeyConcept}
\begin{itemize}
\item The debate (started by \textbf{Ashok Rudra}) tested Green Revolution agriculture against the \textbf{criteria of capitalism}: production for exchange (exchange value over use value), \textbf{wage labour}, monetisation, and \textbf{reinvestment of surplus}. They found no strong positive association --- instead they found the \textbf{emergence of 'gentleman farmers'} (who no longer use their own labour, only supervise).
\item \textbf{Amit Bhaduri} (West Bengal): the mode is \textbf{semi-feudal} --- because of an \textbf{extensive non-legalised sharecropping system}, the \textbf{perpetual indebtedness of small tenants} (landlords are also lenders), \textbf{incomplete access to the rural market}, and \textbf{involuntary exchanges} (the landlord controls the market). Since the power of the feudal lords is not decreasing and they still control marketing, it is semi-feudal (not semi-capitalist).
\item \textbf{Utsa Patnaik} (Punjab/Haryana): the mode is \textbf{pre-capitalistic} --- because capitalism should be judged by \textbf{reinvestment of surplus and accumulation of capital}, not by the free/unfree status of labour (labour is not truly free even under capitalism --- the reserve army of labour). Indian agriculture lacks reinvestment (surplus is spent on rituals, marriages and luxuries), has \textbf{personalised relationships of dependence} between landowner and labourer, and an \textbf{inordinate development of capital in the sphere of exchange} (money-lenders at high rates). It is a transition phase between feudalism and capitalism.
\end{itemize}
\end{KeyConcept}

\begin{QuickRevision}
\begin{itemize}
\item Dialectic = thesis/antithesis/synthesis (Hegel idealist, Marx materialist); dual modernity (Brass); mode of production = means + relations of production.
\item Subaltern school = perspective of the marginalized (tribals' 'irrational' fight was rational); vs Dalit sociology (Kancha Ilaiah English-medium debate).
\item Bhaduri (Bengal): semi-feudal (sharecropping, indebtedness, no market access, involuntary exchange).
\item Patnaik (Punjab): pre-capitalistic (no reinvestment/accumulation; personalised dependence; exchange-sphere capital); gentleman farmers emerging.
\end{itemize}
\end{QuickRevision}

\section{Development, Social Change and Social Transformation}

\begin{KeyConcept}
\begin{itemize}
\item \textbf{Social change} need not alter the structure; \textbf{social transformation} is the \textbf{replacement of problematic social structures} (caste system, class system, unemployment). The Indian state generally pursues social change (incremental), while society demands social transformation --- so \textbf{development itself is often a challenge/obstruction to social transformation}.
\item \textbf{Development is a qualitative concept} (changes in people's lifestyles --- against caste/gender discrimination, conflict, violence), while \textbf{growth is quantitative} (GDP). Types: \textbf{economic development} (equitable growth), \textbf{human development} (literacy, employment and its \emph{quality}), \textbf{social development} (decline in caste/gender violence, discrimination, communal conflict). The idea of development depends on the nature of society (agrarian vs industrial).
\item \textbf{Progress is a value-loaded term}: technological progress (the nuclear bomb) may be social digress; the post-modernists critiqued modernity's 'tomorrow will be better than today'.
\end{itemize}
\end{KeyConcept}

\begin{QuickRevision}
\begin{itemize}
\item Social transformation = replacing problematic structures (not just change).
\item Development = qualitative (lifestyle changes); growth = quantitative; economic/human/social development.
\item Progress = value-loaded (nuclear bomb = tech progress, social digress).
\end{itemize}
\end{QuickRevision}

\section{Dipankar Gupta: From People to Citizens}

\begin{KeyConcept}
\begin{itemize}
\item \textbf{Dipankar Gupta's 'from people to citizens'}: in a multicultural society, people must be \textbf{converted into citizens} to minimise conflict and give minorities a sense of belonging. Citizenship rests on the premise that \textbf{everyone is equal with equal rights over social resources}.
\item So we should demand \textbf{universal health, universal education and universal livelihood} (not 'right to' --- which serves only the unempowered poor, since the empowered already have access; universal systems improve because both empowered and unempowered use them).
\item \textbf{Do not focus much on history}: history brings conflict and false pride (the rewriting of Aurangzeb, Tipu Sultan, the Ramayana); all history is a writer's (biased) interpretation, so we should learn from the \emph{failures and successes} of the past, not individuals.
\item \textbf{Focus less on nationalism and more on democracy}: nationalism shifts discussion to the past (a civilizational 'oneness'), while democracy speaks of the present (equal access, Article 17, caste/gender/communal violence). This makes for a non-ethnic nationalism.
\item \textbf{Make policies for citizens, not for people}: 'people' is a pre-modern, homogeneous category (one religion, one set of values, inequality-friendly), whereas 'citizenship' is modern (different religion/caste/class but equal access and equal voice).
\item \textbf{We need 'citizen elites', not charismatic leaders}: European transformation (the welfare state, the UK's NHS, R\&D centres) was brought by elites at the top; India's successful business/literary/academic people should rally to convert people into citizens. (India's problems: charismatic leaders, class--mass division/crony capitalism, and giving human-social-cultural costs less weight than economic costs.)
\end{itemize}
\end{KeyConcept}

\begin{QuickRevision}
\begin{itemize}
\item Dipankar Gupta: people $\rightarrow$ citizens; universal health/education/livelihood (not 'right to'); history = conflict/false pride (learn from failures, not individuals); democracy over nationalism; policies for citizens not people; 'citizen elites' not charismatic leaders.
\end{itemize}
\end{QuickRevision}

\section{Development-Induced Displacement (Walter Fernandes)}

\begin{KeyConcept}
\begin{itemize}
\item \textbf{Walter Fernandes} (the key scholar of development-induced displacement): development has displaced \textbf{three times} as many people as the Partition of India --- yet this is \textbf{not part of national consciousness}.
\item Estimates: Walter Fernandes --- about \textbf{40 million} displaced; the \textbf{Indian Institute of Public Administration} (54 large dams) --- \textbf{44,182 persons per dam}; \textbf{N. C. Saxena} (former Planning Commission) --- over \textbf{50 million}; the \textbf{World Bank} --- about \textbf{39 million} (3,000+ large dams). (Extremes: \textbf{Arundhati Roy} claims 60 million, \textbf{Surjit Bhalla} 2--3 million.)
\item \textbf{Why the figures are understated}: project proponents \textbf{fudge displacement data} to get government and donor-agency approval (many dams were World Bank-funded).
\item \textbf{Who is displaced}: largely the rural poor --- \textbf{Scheduled Tribes (over 40\% of the displaced, though only $\sim$8\% of the population), Dalits ($\sim$20\%), landless labourers and small/marginal farmers} --- 'Gandhi's last men' --- which is why it never becomes national consciousness.
\item \textbf{Smitu Kothari} documented \textbf{multiple displacement}: people displaced for the Rihand dam, then by a thermal plant, coal mines, railways, industries and a national highway --- five times.
\item \textbf{The spiral of impoverishment}: dehumanisation, disempowerment, psychological trauma, the \textbf{dismantling of the production system} (from forest-dependence to unskilled wage labour), the \textbf{desecration of cultural/religious places}, the \textbf{scattering of kinship groups}, and \textbf{loss of political power} --- leading to the \textbf{production and reproduction of poverty}.
\item \textbf{Arundhati Roy's point}: forest-dwellers who lived off nature (fuel, food, herbs, water) are pushed into a monetised economy at \rupee{}10--20 wages --- so displacement produces poverty. \textbf{Deculturation} leads to \textbf{individualism} and loss of cultural identity. (The World Bank requires both an \textbf{Environmental Impact Assessment and a Social Impact Assessment} before any project.)
\end{itemize}
\end{KeyConcept}

\begin{QuickRevision}
\begin{itemize}
\item Walter Fernandes: development displaced 3$\times$ the Partition; not in national consciousness.
\item Estimates: Fernandes 40m; IIPA 44,182/dam; Saxena 50m; World Bank 39m; Arundhati Roy 60m vs Surjit Bhalla 2--3m.
\item Displaced = ST (>40\%), Dalits ($\sim$20\%), rural poor ('Gandhi's last men'); data fudged for approval.
\item Smitu Kothari: multiple displacement (5$\times$); spiral of impoverishment (dehumanisation, production-system dismantling, deculturation); EIA + SIA.
\end{itemize}
\end{QuickRevision}
